\documentclass[draft]{agujournal2019}
\usepackage{url}
\usepackage{lineno}
\linenumbers

%% Journal name — change to match target journal
%% Options: JGR Atmospheres, JGR Oceans, GRL, GC Cubed, etc.
\journalname{JGR: Atmospheres}

\begin{document}

\title{ENSO-conditioned evolution of global mean surface temperature}

\authors{{{AUTHOR NAME}}\affil{1}}

\affiliation{1}{Department, Institution, City, Country}

\correspondingauthor{{{AUTHOR NAME}}}{email@institution.edu}

\begin{keypoints}
\item Key point 1 here (under 140 characters)
\item Key point 2 here (under 140 characters)
\item Key point 3 here (under 140 characters)
\end{keypoints}

\begin{abstract}
ENSO forecasts are skillful, and the influence of ENSO on global
  mean surface temperature (GMST) is well established. These facts
  suggest that future GMST behavior can be anticipated based on future
  expected ENSO conditions. Here we examined how well the June--May trajectory of monthly
  GMST can be anticipated from information available near the beginning of June. The information we used were recent GMST evolution and the
  Nino-3.4 values expected during the upcoming boreal winter. Principal component analysis and dimension reduction of both GMST and
  Nino-3.4 data simplified the problem greatly. Further simplifications based on the empirical orthogonal
  function structure led to a simple, interpretable model in which the
  June--May trajectory of monthly GMST is conditioned on the average
  of prior 12 months and the expected December value of Nino-3.4. We found that an expected December Nino-3.4 anomaly of
  1\textdegree C is associated a GMST increase from the value of the
  previous year of about 0.025\textdegree C in boreal summer and a
  peak increase of about 0.1\textdegree C in the following February.
\end{abstract}

\section*{Plain Language Summary}
Changes in global mean surface temperature (GMST) changes are due to
  both anthropogenic and natural forcing as well as natural
  variability, with El Niño-Southern Oscillation (ENSO) being the
  dominant contributor from natural variability. In this study, we
  examined how much of the coming June-to-May evolution of global mean
  temperature can be anticipated from two pieces of information
  available in advance: the recent global temperature level and the
  expected state of ENSO over the same period. We found in a simple
  model that an expected December Nino-3.4 anomaly of 1\textdegree C
  is associated a GMST increase of about 0.025\textdegree C in boreal
  summer and a peak increase of about 0.1\textdegree C in the
  following February.

\section{Introduction}

Global mean surface temperature (GMST) is trending upward due
    to human-caused increases in greenhouse gases. Superimposed on
    this upward trend, GMST also varies in response to natural forcing
    (e.g., volcanoes) and internal variability. ENSO is the dominant component of internal variability on
    seasonal timescales, and its influence on GMST is well established
    \cite<e.g.,>[]{Trenberth2002evolution}. During warm ENSO
    conditions (El Ni\~no) GMST tends to be elevated relative to the
    long-term trend, while during cool ENSO conditions (La Ni\~na),
    the opposite tends to occur. \citeA{Lean2008natural} attributed about 0.23\textdegree C of
    GMST warming to the 1997 ``super'' El Ni\~no in June to November
    of that year. The World Meteorological Organization recently noted ``The
    most recent El Ni\~no, in 2023-24, was one of the five strongest
    on record and it played a role in the record global temperatures
    we saw in 2024'' \cite{WMO2026enso}.

  
ENSO is predictable several seasons in advance. Each month
    meteorological and climate forecast centers around the world issue
    operational ENSO forecasts that extend typically two to four
    seasons into the future.  ENSO forecasts issued after the boreal
    spring predictability barrier are especially skillful
    \cite<e.g.,>[]{Heureux:Chapter}. In early 2026, forecasts of a strong El Nino have led to
    speculation in the press that it could boost GMST by about
    0.2\textdegree C \cite<e.g.,>[]{KeithLucas2026SuperElNino}.
    While this estimate appears plausible, exactly how a given
    forecast of ENSO conditions translates into an expected GMST
    anomaly is less clear, both in terms of timing and amplitude. For instance, it remains difficult to say, in a quantitative
    and principled manner, what GMST evolution is expected over the
    next 12 months given its recent values and the ENSO conditions
    forecast for this winter. One reason for this difficulty is that previous analyses
    either pooled all months together \cite{Trenberth2002evolution}
    or focused on the ENSO–GMST relationship only at the lag at which
    the correlation is strongest
    \cite{Lean2008natural,Foster2011global}.

  
  
Here, we address this gap and present a simple and
    interpretable methodology for anticipating GMST evolution at
    monthly resolution over the June--May 12-month period, given a
    forecast of ENSO conditions over the same period.  We chose this
    period because ENSO forecasts initialized after boreal spring
    typically have useful skill through the following winter. We conditioned our analysis on observed ENSO evolution and
    acknowledge imperfect ENSO forecasts as an additional source of
    uncertainty, analogous to the perfect-prog framework in weather
    output statistics \cite{Glahn72,Klein1959objective}. Data sources and statistical methods are described in Section
    \ref{sec:data-methods}.  Section \ref{sec:results} introduces our
    data-compression strategies and a sequence of regression models,
    progressing from general formulations to simpler ones that are
    more interpretable and nearly as skillful. We also show why these
    successive model reductions are physically and statistically well
    motivated, and evaluate model performance using hindcasts and
    summary skill measures. Conclusions and discussion are given in
    Section \ref{sec:summary}.

\section{Data and Methods}\label{sec:data-methods}

Our analysis used monthly data from June 1979 through May
    2025, comprising 46 June--May ENSO years. Global mean surface temperature (GMST) was computed as the
    area-weighted global mean of gridded surface temperature from NOAA
    GlobalTemp v6.1.0 \cite{Huang2022improvements}. NOAA GlobalTemp
    uses the current WMO climatological standard normals period,
    1991--2020, as the reference period for temperature anomalies. Its
    1850--1900 mean is $-0.77^\circ$C, and this value was subtracted
    to form the preindustrial anomalies shown in the figures. The Niño-3.4 index was computed from ERSSTv5
    \cite{Huang:ERSST5} as the area-weighted sea surface temperature
    over 5\textdegree S--5\textdegree N, 170\textdegree
    W--120\textdegree W.  Anomalies are shown relative to the
    1991--2020 base period.

Principal component analysis (PCA) was applied to the
    $12 \times 46$ Niño-3.4 and GMST data matrices after removal of
    monthly means (row averages). In the data matrices, rows
    correspond to months and columns to years. Double-centered PCA was applied to the same $12 \times 46$
    Niño-3.4 and GMST data matrices after removal of both monthly
    means (row averages) and June--May means (column averages). We followed the PCA convention in which principal components
    (PCs) have unit variance and empirical orthogonal functions (EOFs)
    carry the physical units \cite{DelsoleTippettBook}. Data were reconstructed as the product of EOFs and PCs, with
  the removed means added back. Coefficient of determination ($R^2$), root-mean-square error
  (RMSE), and mean-squared-error skill score (MSESS) were used as
  summary measures of regression model performance. Statistical significance of differences in RMSE was assessed
    using the Wilcoxon signed-rank test applied to paired squared
    errors. Because the RMSE values were computed using the same GMST
    data, RMSE values for different models are not independent, and an
    $F$-test for equality of variances is inappropriate
    \cite{TDS:Comparing:2014}. The Wilcoxon signed-rank test
    effectively tests whether the median difference squared errors is
    zero. Statistical significance for RMSE and MSESS is identical
    because MSESS is a monotonic function of RMSE.

\section{Results}\label{sec:results}

A regression model that included all plausible predictors and
    predictands would be high-dimensional: 24 predictors and 12
    predictands. The predictors would be the previous 12 months of
    GMST together with the next 12 months of Niño-3.4, while the
    predictands would the next 12 months of GMST. Such a model would
    be prone to overfitting and difficult to interpret. Therefore, as
    a first step, we applied principal component analysis (PCA) to the
    data. PCA provides an unsupervised and systematic dimension
    reduction that is based on explained variance and ignores
    relationships between predictors and predictands. Previous work has shown that 12-month Nino-3.4 trajectories
    are well approximated by one or two PCs, both in observations and
    in seasonal forecasts \cite{Bunge2009verified,Tippett2019:ENSO}.
    Consistent with these results, we found that approximately 95\% of
    the variability of June--May Nino-3.4 segments is explained by the
    first two PCs (Fig.\ \ref{fig:EOF-variance-explained}a). Moreover, we found that the same is true of
    GMST---approximately 95\% of the variability of a June--May GMST
    trajectory is explained by the first two PCs (Fig.\
    \ref{fig:EOF-variance-explained}b).

  
We now examine the EOFs that go with those PCS. EOF1 of Niño-3.4 peaks in December, reflecting the seasonal
    phase locking of ENSO (Fig.\ \ref{fig:fig1}a).  The correlation of
    PC1 of Nino-3.4 with the December Nino-3.4 index exceeds 0.98
    (compare PC1 of Nino-3.4 with the standardized value of December
    Nino-3.4 in Fig.\ \ref{fig:fig1}c).  EOF2 of Niño-3.4 is a dipole
    in time centered on December and can account for shifts of the
    Nino-3.4 peak away from December. EOF1 of GMST is nearly uniform across months (Fig.\
    \ref{fig:fig1}b), while PC1 of GMST shows a clear upward trend
    (Fig.\ \ref{fig:fig1}d). This behavior indicates PC1 represents
    the uniform (across calendar months) upward trend of GMST. EOF2 of GMST is nearly uniform and negative June--November
    until it rises in December and peaks in January and February,
    while PC2 shows no visible trends. This behavior suggested that
    PC2 might represent the well-known delayed response of GMST to
    ENSO via the atmospheric bridge
    \cite{Alexander2002atmospheric,Klein1999remote}. Reconstructions of Nino-3.4 and GMST show that 2 EOFs are
    adequate to capture interannual variability ($R^2$ values of about
    95\%; Fig.\ \ref{fig:fig1}e, f), although more unresolved
    subannual variability remains visible in the GMST time series than
    in the Niño-3.4 time series. The two-PC reconstruction of GMST
    sets an upper bound on the possible accuracy of a two-PC model.

\begin{figure}
  \noindent \includegraphics[width=\textwidth]{/Users/tippett/claude/enso-gmt-prediction/plots/eofs_pcs.pdf}
  \noindent \includegraphics[width=\textwidth]{/Users/tippett/claude/enso-gmt-prediction/plots/reconstruction_2pc.pdf}
  \caption{The leading two EOFs of June--May segments of
      (a) Nino-3.4 and (b) global mean surface temperature (GMST). The
      corresponding leading two PCs of (c) Nino-3.4 and (d) GMST.  The
      standardized December value of Nino-3.4 is also shown in panel
      (c) for comparison with PC1. Two-PC reconstructions of (e)
      Nino-3.4 and (f) GMST with root mean squared error (RMSE) and
      R-squared values in the title.}
  \label{fig:fig1}
\end{figure}

We next examined the cross correlations and lag correlations of
    the first two PCs of Niño-3.4 and GMST. The strongest simultaneous correlation is the 0.61 correlation
    between N34-PC1 and GMST-PC2, which supports our hypothesis that
    GMST-EOF2 represents the delayed GMST response to ENSO
    (Fig. \ref{fig:fig2}a). The other simultaneous cross-correlations are
    weak ($R^2 < 4\%$). The strongest lagged relation between PCs is the 0.92
    correlation between GMST-PC1$(t)$ and GMST-PC1$(t+1)$ which
    reflects the upward trend in GMST-PC1 and the resulting strong
    year-to-year persistence (Fig. \ref{fig:fig2}b).  In contrast, the
    lag-1 autocorrelation of GMST-PC2 is small and negative
    ($r=-0.18$), consistent with the interpretation of GMST-PC2 as an
    ENSO response that does not persist from one year to the next. The source of the correlation of N34-PC2$(t)$ with
    N34-PC1$(t+1)$ is less obvious, though it is visually evident in
    the PC time series for both warm and cool ENSO events (Fig.\
    \ref{fig:fig1}c). \textbf{Thoughts?}. The correlation between N34-PC2$(t)$ and GMST-PC2$(t+1)$
    ($r=0.37$) is explained almost entirely by the pathway
    N34-PC2$(t) \rightarrow$ N34-PC1$(t+1) \rightarrow$
    GMST-PC2$(t+1)$. After accounting for N34-PC1$(t+1)$, the partial
    correlation between N34-PC2$(t)$ and GMST-PC2$(t+1)$ is only 0.04. The correlation analysis above suggests that little is lost by
    using two predictors, GMST-PC1$(t)$ and N34-PC1$(t+1)$, instead of
    four, GMST-PC1$(t)$, GMST-PC2$(t)$, N34-PC1$(t+1)$, and
    N34-PC2$(t+1)$. Indeed, the fit obtained with these two predictors
    is nearly identical to that obtained with all four predictors
    (compare Figs.\ \ref{fig:fig2}c and
    \ref{fig:fit_full_recon}). Moreover, in the four-predictor model,
    only the coefficients of GMST-PC1$(t)$ and N34-PC1$(t+1)$ are
    large and statistically significant (Table \ref{tab:full_model}).

\begin{figure}
  \noindent \includegraphics[width=\textwidth]{/Users/tippett/claude/enso-gmt-prediction/plots/pc_correlations.pdf}
  \noindent \includegraphics[width=\textwidth]{/Users/tippett/claude/enso-gmt-prediction/plots/fit_reduction1.pdf}
  \caption{Panel (a) cross correlations (simultaneous)
      between the first two PCs of Nino-3.4 and and global mean
      surface temperature (GMST). Panel (b) cross correlations
      (one-year lag) between the first two PCs of Nino-3.4 and
      GMST. Panel (c) shows the fit of GMST time series by the two-PC
      model (see text for details). Dots show the start of June--May
      segments. The RMSE and R-squared of the four-PC model is shown
      in the title for comparison.}
  \label{fig:fig2}
\end{figure}

While the two-PC predictor model is both effective and
    parsimonious, there are further simplifications that can improve
    interpretability. In particular, the predictors GMST-PC1$(t-1)$ and N34-PC1$(t)$
    are defined through the PCA and therefore depend on both the
    specific data used and the details of the PCA
    calculation. Replacing them with quantities that closely
    approximate them, but whose calculation is more transparent, would
    yield a more explainable model. In the first simplification, we replaced N34-PC1 with the
    December Niño-3.4 value, which correlates with N34-PC1 at $r=0.99$
    (Fig.\ \ref{fig:fig1}c). This close correspondence follows from
    the EOF structure: EOF1 peaks in December, while EOF2 is close to
    zero in December. The resulting model performs essentially as well as the two-PC
    predictor model (Fig.\ \ref{fig:fit_r2_recon}) while being easier
    to interpret. In particular, the dependence of GMST on ENSO can
    now be described without reference to PCs, for example: ``given
    December Niño-3.4 $=x$, the anticipated June--May GMST trajectory
    is \ldots''. For the second simplification, we replaced GMST-PC1. Recall
    that GMST-EOF1 is nearly uniform across months, which suggests a
    decomposition in which the first mode is simply the June--May mean
    GMST (which we call the level), while the second mode captures as
    much as possible of the remaining month-to-month structure and is
    therefore expected to resemble GMST-PC2. We implemented this idea by decomposing each June--May GMST
    trajectory into its June--May level together with a PCA of the
    departures from that mean. Because this decomposition removes both
    monthly means and June--May segment means, we call it
    double-centered PCA. Double-centered PCA is not variance-optimal so its modes
    explain less variance, but the loss is small: the variance
    explained by the first two modes decreases from 94.8\% to 94.3\%
    (Fig.\ \ref{fig:EOF-variance-explained-dc}). The leading EOF from
    the double-centered decomposition is similar to GMST-EOF2 (Fig.\
    \ref{fig:fig3}a), and the corresponding PC is also similar to
    GMST-PC2, with correlation 0.993 (Fig.\ \ref{fig:fig3}b). We call the model with these two simplifications (N34-PC1
    replaced by December Nino-3.4 and GMST-PC1 replaced by level), the
    simplified model. The performance of the  simplified model is nearly
    identical to that of the two-PC predictor model (Fig.\
    \ref{fig:fig3}c). The simplified model improves interpretability by enabling
    statements of the form: ``given recent GMST level $\bar{g}$ and
    December Niño-3.4 $=x$, the June--May GMST trajectory is \ldots''.

  
A remaining question is how much is gained by conditioning on
    ENSO. Because of the strong upward trend in GMST,
    persistence-based predictions are already a nontrivial
    benchmark. We therefore compared the simplified model with two
    persistence baselines: one based on the preceding June--May mean 
    and one based on the preceding May GMST value. The simplified model has lower RMSE than the June--May-mean
    baseline in every month and lower RMSE than the May baseline from
    August through May (Fig.\ \ref{fig:fig3}d). Its clearest and most
    statistically robust advantage occurs from December through May,
    when the delayed GMST response to ENSO is expected to be
    strongest. The corresponding MSE skill scores show that the simplified
    model improves substantially on both persistence baselines, with
    the largest gains generally occurring during boreal winter and
    spring (Fig.\ \ref{fig:fig3}e).

\begin{figure}
  \noindent \includegraphics[width=\textwidth]{/Users/tippett/claude/enso-gmt-prediction/plots/eofs_pcs_dc.pdf}
  \noindent \includegraphics[width=\textwidth]{/Users/tippett/claude/enso-gmt-prediction/plots/fit_idea2.pdf}
  \noindent \includegraphics[width=\textwidth]{/Users/tippett/claude/enso-gmt-prediction/plots/rmse_by_month.pdf}
  \caption{Comparison of the standard (a) EOFS and (b) PCs
      of global mean surface temperature (GMST) with double-centered
      ones. Panel (c) shows the fit of GMST time series by the
      simplified model (see text for details). (d) RMSE and MSESS of
      the simplified model compared to persistence
      baselines.}
  \label{fig:fig3}
\end{figure}

Summary performance metrics are informative, but they do not
  show what the simplified model looks in individual years.
  To address this point, Figs.\ \ref{fig:fig4}a--f shows
  representative hindcasts for several recent years together with
  an illustrative conditional forecast for 2026--27. The prediction intervals are fairly wide, but in most cases
  they provide a reasonable characterization of uncertainty. In the cool ENSO years shown (2007, 2010, and 2020), both
  the simplified model and the observed GMST generally lie below
  the level of the previous 12 months, with a pronounced dip
  around January and a partial rebound by May. The amplitude of
  the simplified model is smaller than that of the observed GMST,
  as expected from a regression model; in particular, the
  relevant coefficient of determination is only $R^2=0.40$
  (Table \ref{tab:model-sm}). In the warm ENSO cases, the behavior is broadly opposite,
  with elevated GMST during boreal winter and early spring.
  An important exception is 2023--24, when the sharp warming in
  September is not consistent with the expected timing of the
  delayed ENSO influence and lies outside the prediction
  intervals. An attractive feature of the simplified model is that, once
  the recent GMST level is specified, the month-to-month shape of
  the predicted GMST response depends only on December
  Niño-3.4. In particular, the seasonal pattern is fixed and only
  its amplitude changes with the December Niño-3.4 value. This structure makes it straightforward to describe the GMST
  response to a specified ENSO condition. Figure \ref{fig:fig4}g
  shows the implied June--May GMST response per
  $+1^\circ$C of December Niño-3.4. The response is about
  $0.025^\circ$C during boreal summer, rises sharply beginning in
  December, and peaks at slightly more than $0.10^\circ$C in
  February. The same analysis can be carried out using relative
  Niño-3.4, or RONI \cite{LHeureux2024relative}. The implied
  GMST response is very similar, although slightly weaker in
  amplitude, consistent with relative Niño-3.4 being somewhat
  smaller than standard Niño-3.4 during warm ENSO conditions. We also show an illustrative 2026--27 GMST trajectory based
  on the currently observed June 2025--February 2026 mean of
  1.26$^\circ$C above 1850--1900 together with an assumed
  December 2026 Niño-3.4 value of $+2.0^\circ$C. Under these conditions, the predicted trajectory peaks at
  about 1.51$^\circ$C above preindustrial in February 2027,
  close to the expected seasonal peak of ENSO influence.

\begin{figure}
  \noindent \includegraphics[width=\textwidth]{/Users/tippett/claude/enso-gmt-prediction/plots/hindcast_panels.pdf}
  \noindent \includegraphics[width=\textwidth]{/Users/tippett/claude/enso-gmt-prediction/plots/forecast.pdf}
  \caption{Observed June--May global mean surface
      temperature (GMST) and simple model fits for (a) 2007-8, (b)
      2009-10, (c) 2010-11, (d) 2015-16, (e) 2020-21, and (f)
      2023-24. Panel (g) shows a 2026-27 illustration for Dec 2026
      Nino-3.4 of 2.0\textdegree C.}
  \label{fig:fig4}
\end{figure}

\section{Summary and discussion}\label{sec:summary}

Summarize the main findings in reverse order. Start with the
    ending and explain how we got there. Implications. 2023-24. How the work here differs from the fitting approaches of
    \citeA{Lean2008natural} and \citeA{Foster2011global}. A paragraph about prediction of GMST with seasonal climate
    models and the problems that \citeA{Tippett2024trends} found
    regarding multimodel too over dispersed and initialization
    etc. https://agupubs.onlinelibrary.wiley.com/doi/full/10.1029/2024GL110703. Future work: short lead forecasts should be more accurate. Is
    this lost by doing the entire trajectory at once? Other ideas?

Here we addressed the question of how much of the year-to-year
  evolution of monthly GMST can be anticipated from
  information available in advance, and in particular from the current
  level of GMST and the expected state of ENSO. We end with a simple practical statement: given the recent
  June--May mean GMST and the expected December
  Niño-3.4 value, we can anticipate the full June--May evolution of
  GMST. In that formulation, the recent GMST sets the
  overall level of the next year's trajectory, while December
  Niño-3.4 controls much of the within-year wintertime amplification. The resulting two-predictor model is both interpretable and
  effective: it reproduces the main year-to-year evolution of global
  mean temperature with skill nearly identical to less interpretable
  PC-based formulations. That simple formulation was reached by first noting that the
  leading EOF of GMST is nearly uniform across
  months, so it is well represented by the June--May mean rather than
  by a less transparent principal component. Likewise, the leading EOF of Niño-3.4 peaks in December, and its
  leading PC is therefore almost equivalent to the December Niño-3.4
  value itself. These simplifications worked because the June--May trajectories
  of both Niño-3.4 and GMST are intrinsically
  low-dimensional: two modes explain about 95\% of the variance in
  each case. The cross-correlation structure then showed why a parsimonious
  predictor set is sufficient: persistence in GMST
  level carries information from one year to the next, while ENSO
  primarily projects onto the delayed wintertime shape of the global
  mean temperature trajectory. In physical terms, the analysis quantifies a familiar picture:
  long-term warming sets the background state, and ENSO modulates the
  timing and amplitude of departures from that background, especially
  in boreal winter. Thus, much of the apparent complexity of monthly GMST evolution
  can be reduced to a slowly varying level plus an ENSO-conditioned
  seasonal shape.

Unlike \citeA{Lean2008natural} and \citeA{Foster2011global},
  which regress temperature on external explanatory indices, the
  present approach is explicitly autoregressive: the recent GMST level
  itself is the leading predictor of the next year's evolution. The most important difference is therefore that their fitting
  frameworks do not use the recent temperature state itself as a
  predictor, whereas here the prior June--May mean GMST carries
  forward the slowly varying background state and sets the baseline
  for the upcoming June--May trajectory. ENSO then enters as a second, physically interpretable predictor,
  not simply as an explanatory covariate, but as a predictor of the
  within-year departure from that baseline. In particular, our analysis exploits the seasonality of the ENSO
  response: the effect of ENSO on GMST is not
  uniform through the year, but is strongest in boreal winter, so the
  expected December Niño-3.4 value helps determine the seasonal shape
  of the upcoming temperature trajectory. Thus, the model separates two distinct sources of predictability:
  persistence of the recent global temperature level and the seasonal,
  lagged imprint of ENSO. A further difference is structural: rather than fitting
  month-by-month temperature anomalies in calendar time, we model the
  June--May trajectory as a whole and exploit its low-dimensional
  structure. This yields a compact two-predictor forecast model with a clear
  interpretation: the recent temperature mean sets the level, and the
  expected December Niño-3.4 value shapes the seasonal departure from
  that level, especially in boreal winter. In that sense, \citeA{Lean2008natural} and
  \citeA{Foster2011global} are primarily explanatory or adjustment
  frameworks, whereas the present work is an explicitly predictive,
  autoregressive framework for anticipating the next year's monthly
  GMST evolution.

Initialized seasonal climate models provide a physically
  self-consistent way to predict GMST, because they
  evolve the coupled climate system forward rather than fitting global
  mean temperature statistically after the fact. However, \citeA{Tippett2024trends} showed that current systems
  still have important shortcomings for this problem, even though they
  do exhibit skill beyond a simple trend baseline. In particular, much of that skill is tied to ENSO, but the
  forecast quality depends strongly on season, forecast uncertainty is
  sensitive to initialization, individual models tend to be
  underdispersed at short lead, and the multimodel ensemble tends to
  be overdispersed. The multimodel mean captures much of the broad evolution of
  recent GMST, but it underestimated the record
  warmth of 2023. Thus, seasonal climate models are powerful and physically
  complete, but they are not yet clean, plug-and-play tools for
  monthly GMST prediction. There is a useful analogy here with the older two-tier seasonal
  prediction strategy used widely in the 1990s and 2000s, in which SST
  was predicted first and the atmospheric response was then computed
  separately with an AGCM, largely because coupled model biases were
  too large to rely on a fully coupled forecast alone. Our approach is similar in spirit: rather than relying on a full
  coupled prediction system to get every part of the problem right at
  once, we isolate the most predictable large-scale components. In particular, we use the slowly varying GMST
  background state together with the seasonally dependent ENSO
  influence in a simpler predictive framework.

Extend the analysis from a proof-of-concept predictive model to
  a fully out-of-sample forecast framework with cross-validation and
  real-time prediction experiments. Quantify forecast uncertainty more carefully, including
  out-of-sample prediction intervals and calibration of probabilistic
  forecasts. Examine whether additional predictors beyond the prior June--May
  mean GMST and December Niño-3.4 can add useful
  skill, especially for years with unusual evolution. In particular, it will be useful to test whether volcanic
  forcing, solar variability, or other modes of internal variability
  provide incremental predictive information beyond the present
  two-predictor model. Explore alternative representations of ENSO, including indices
  that distinguish event amplitude, longitude, or temporal evolution,
  to determine whether they better capture the seasonally varying
  GMST response. Extend the framework to initialized ENSO forecasts, so that the
  December Niño-3.4 predictor is itself predicted operationally rather
  than prescribed. Compare the statistical framework more directly with initialized
  seasonal climate model forecasts of GMST, with
  particular attention to reliability, dispersion, and sensitivity to
  initialization. Investigate whether the same low-dimensional trajectory approach
  can be applied to regional or zonal mean temperatures, where the
  forced signal is weaker and the seasonal ENSO response may be more
  structured. More broadly, future work should clarify how far monthly global
  mean temperature predictability can be pushed using transparent,
  physically interpretable predictors before the added complexity of
  larger forecast systems becomes necessary.

NOAAGlobalTemp version 6.1 data are from \url{https://www.ncei.noaa.gov/data/noaa-global-surface-temperature/v6/access/gridded/NOAAGlobalTemp_v6.1.0_gridded_s185001_e202602_c20260308T114550.nc}. NOAA Extended Reconstructed SST V5 (ERSSTv5) data are provided
  by the NOAA PSL, Boulder, Colorado, USA, from their website
  \url{https://psl.noaa.gov}.

Please include a comprehensive conflict of interest statement that reflect all conflicts of interest for all involved authors. If there are no conflicts of interest please state, “The authors declare there are no conflicts of interest for this manuscript.”

\clearpage

\section*{Open Research Section}
Data and code availability statement here. Data cannot be listed as
``available from the authors.''

\section*{Conflict of Interest}
The authors declare no conflicts of interest.

\acknowledgments
Enter acknowledgments here.

\bibliography{all}

\section{Supplemental/not shown}

\renewcommand\thefigure{S\arabic{figure}}

\setcounter{figure}{0}

\renewcommand\thetable{S\arabic{table}}

\setcounter{table}{0}

\begin{figure}
  \noindent \includegraphics[width=\textwidth]{/Users/tippett/claude/enso-gmt-prediction/plots/variance_explained.pdf}
  \caption{Variance per mode and cumulative explained
      variance of PCA applied to (a, c) Nino-3.4 and (b, d) global
      mean temperature (GMST). The dashed lines in panels a, b indicate
      the 5\% level and in panels c, d the 90\%
      level.}
  \label{fig:EOF-variance-explained}
\end{figure}

\begin{figure}
  \noindent \includegraphics[width=\textwidth]{/Users/tippett/claude/enso-gmt-prediction/plots/fit_full_model.pdf}
  \caption{Reconstruction fit of the full two-PC
      model.}
  \label{fig:fit_full_recon}
\end{figure}

    
Predictors: GMST PC1, PC2 at time $t$
            and Niño-3.4 PC1, PC2 at time $t+1$. In-sample fit: RMSE$=0.11$°C, $R^2=0.86$. Gap to 2-PC ceiling ($R^2=0.95$) reflects
            year-to-year variability not explained
            by the chosen predictors. This full model is the benchmark;
            reductions next.

\begin{figure}
  \noindent \includegraphics[width=\textwidth]{/Users/tippett/claude/enso-gmt-prediction/plots/fit_reduction2.pdf}
  \caption{Reconstruction fit of reduction 2.}
  \label{fig:fit_r2_recon}
\end{figure}

    
Replace Niño-3.4 PC1 with the raw \textbf{December}
            Niño-3.4 value ($r=0.99$ with PC1). December is where EOF1 peaks and EOF2 $\approx 0$
            --- a single month proxies the full trajectory. RMSE$=0.114$, $R^2=0.85$ --- same skill as
            Reduction~1, zero loss vs the PC formulation. Enables direct physical statements:
            ``given December Niño-3.4 $= x$, the
            June--May GMST trajectory is $\ldots$''.

\begin{figure}
  \noindent \includegraphics[width=\textwidth]{/Users/tippett/claude/enso-gmt-prediction/plots/variance_explained_dc.pdf}
  \caption{Variance per mode and cumulative
      variance for double-centered PCA.}
  \label{fig:EOF-variance-explained-dc}
\end{figure}

    
Remove annual cycle \emph{and} segment mean
            (Jun--May average) before PCA. Segment mean captures \textbf{90.9\%} of
            total variance — this is the ``level.''. Remaining PCA modes capture the
            within-year ``shape'' ($\sim$5\%). Segment mean $\equiv$ GMST-PC1 (correlation
            $r = 1.00$ since EOF1 is nearly flat).

\begin{figure}
  \noindent \includegraphics[width=\textwidth]{/Users/tippett/claude/enso-gmt-prediction/plots/corr_dec_n34_by_month.pdf}
\end{figure}

Expected: positive correlations peaking
            near December (ENSO peak season) ---
            confirmed. Unexpected: Jun--Aug correlations are
            equally strong but \emph{negative}. The small Jun--Aug regression coefficients
            (see left panel of previous figure) are
            not sampling noise --- they represent
            genuine signal. Hypothesis: ENSO events tend to alternate
            in sign year to year.  When Dec Niño-3.4
            is warm (El Niño), the preceding Jun--Aug
            was more likely cool (residual La Niña
            from the prior year), producing negative
            correlation. Implication: the DC residual in early months
            carries memory of the \emph{prior} ENSO
            event, not captured by Dec Niño-3.4 alone.

\begin{figure}
  \noindent \includegraphics[width=\textwidth]{/Users/tippett/claude/enso-gmt-prediction/plots/pc_correlations_dc_pred.pdf}
  \caption{Cross correlations for idea 2.}
  \label{fig:idea-2-cross-corr}
\end{figure}

    
GMST mean$(t)$ predicts GMST mean$(t+1)$
            strongly ($r \approx 0.92$) --- persistence
            dominates the level. Dec Niño-3.4$(t+1)$ predicts DC-PC1$(t+1)$
            ($r \approx 0.63$) --- ENSO drives the
            within-year shape. DC-PC1$(t)$ carries little additional
            information beyond the level. Two predictors, two predictands: a
            minimal, interpretable system.

\begin{figure}
  \noindent \includegraphics[width=\textwidth]{/Users/tippett/claude/enso-gmt-prediction/plots/pc_correlations_dc.pdf}
\end{figure}

GMST mean$(t)$ predicts GMST mean$(t+1)$
            strongly ($r=0.92$, persistence) and
            new PC2$(t+1)$ weakly ($r=0.22$). Dec Niño-3.4$(t+1)$ predicts new
            PC1$(t+1)$ strongly ($r=0.63$) and
            GMST mean$(t+1)$ weakly ($r=0.17$). new PC2$(t+1)$ correlates with
            nothing strongly --- likely noise
            (new EOF2 $\approx$ flat, $R^2=0.12$).

\begin{table}
  \input{/Users/tippett/claude/enso-gmt-prediction/tex/regression_table_full.tex}
\caption{Regression table for the four-PC model}\label{tab:full_model}
\end{table}

\begin{table}
  
  \input{/Users/tippett/claude/enso-gmt-prediction/tex/regression_table_r1.tex}
  \caption{Regression table for the two-PC model}\label{tab:model-r1}
  \end{table}

\begin{table}
  
  \input{/Users/tippett/claude/enso-gmt-prediction/tex/regression_table_r2.tex}
  \caption{Regression table for N34-PC1 replaced by Dec Nino-3.4 in the two-PC model.}
\end{table}

\begin{table}
  
  \input{/Users/tippett/claude/enso-gmt-prediction/tex/regression_table_dc2.tex}
  \caption{Regression table for the simplified model.}
  \label{tab:model-sm}
\end{table}

\end{document}
